\subsection{A complementary factor and the remaining arithmetic question at thirty-one}
\label{subsec:thirtyone-gate}

The next calculation concerns one cyclic quotient of the Fermat curve.
It identifies its only possible rational rank-zero factor, but does not
establish a new section classification. Let $C/\Q$ be the smooth proper
model of
\[
 y^{31}=x(1-x)^5,
\]
put $J=\Jac(C)$, and define
\[
 c(x,y)=\left(\frac1{1-x},-\frac{y^5}{1-x}\right),\quad
 B=C/\langle c\rangle,\quad A=\Jac(B),\quad
 e=\frac{1+c+c^2}{3}.
\]
Let $P$ be the image of $1-e$ up to rational isogeny; equivalently,
one may take the connected image of $2-c-c^2$. Put
\[
 K=\Q(\zeta_{31}),\quad H=\langle5\rangle=\{1,5,25\},\quad
 M=K^H,\quad F_0=M^+.
\]

\begin{proposition}[The rational factors and their rank bounds]
\label{prop:thirtyone-rational-factors}
There is a rational isogeny $J\sim_\Q A\times P$, where
$\dim A=5$, $\dim P=10$, and
\[
 \operatorname{End}^0_\Q(J)=\Q\times\Q(\zeta_3),\qquad
 \operatorname{End}^0_\Q(A)=\Q,\qquad
 \operatorname{End}^0_\Q(P)=\Q(\zeta_3).
\]
Both factors are $\Q$-simple, and
\[
 \operatorname{rank}A(\Q)\ge1,\qquad
 \operatorname{rank}J(\Q)\le4,\qquad
 \operatorname{rank}P(\Q)\in\{0,2\}.
\]
Thus $P$ is the only possible positive-dimensional rational
rank-zero isogeny factor of this particular $J$.
\end{proposition}
\begin{proof}
The degree-thirty-one map to the $x$-line has three fully ramified
points, so Riemann--Hurwitz gives $g(C)=15$.
Direct substitution gives
\[
 c^2(x,y)=\left(\frac{x-1}{x},-\frac{y^{25}}{x(1-x)^4}\right),
 \qquad c^3=1.
\]
The fixed $x$-coordinates satisfy $x^2-x+1=0$.
For a primitive sixth root $\rho$ with $\rho+\rho^{-1}=1$,
the point $(\rho,\rho^2)$ is fixed. Multiplication of its
$y$-coordinate by $\zeta_{31}^i$ changes the fixed-point condition
to $5i=i$, so this is the unique fixed point in that fiber.
The conjugate fiber contributes one more, and the three branch
points are cyclically permuted. Hence
$28=3(2g(B)-2)+4$, giving $g(B)=5$ and the stated dimensions.
Norm and pullback give the displayed rational isogeny.

The Jacobi-sum CM type on $K$ is
\[
 S=\{1,2,3,4,5,7,8,9,10,13,14,15,19,20,25\};
\]
it consists of the residues $b$ with $b+(5b\bmod31)<31$.
Its multiplicative stabilizer is exactly $H$. In the cyclic
quotient of order ten, with generator induced by $3\bmod31$,
the type is
\[
 T=\{0,1,2,4,8\}\subset\mathbb Z/10\mathbb Z,
\]
whose translation stabilizer is trivial. The CM decomposition
of \citet[Lemma~1.6, pp.~206--207]{GrossRohrlich1978} therefore
gives $J_{\overline\Q}\sim D^3$, where $D$ is absolutely simple
of dimension five with CM by $M$. In particular
$\operatorname{End}^0_{\overline\Q}(J)=M_3(M)$.
The automorphism $c$ and the cyclotomic endomorphisms satisfy
$cac^{-1}=\sigma_5(a)$. Their crossed product $K\rtimes H$ is
$\operatorname{End}_M(K)=M_3(M)$, so these endomorphisms generate
the whole geometric algebra. Galois acts on the unique expansion
$\sum_{j=0}^2a_jc^j$ through its usual action on $a_j\in K$,
and fixes $c$. The rational fixed algebra is therefore
$\Q[c]=\Q\times\Q(\zeta_3)$. Its two central factors give the
claimed endomorphism algebras and rational simplicity.
The same argument over $F_0$ gives $F_0[c]$; in particular $A$
is absolutely simple and $A_{F_0}$ has real multiplication by $F_0$.

\citet[Theorem~2.1, pp.~208--209]{GrossRohrlich1978} gives a
rational divisor class of infinite order on this cyclic quotient:
\[
 D_0=[(\rho,\rho^2)+(\rho^{-1},\rho^{-2})-2\infty].
\]
The exclusions in that theorem are $s=1,15,29$, so $s=5$ is
permitted. The two displayed points are fixed by $c$.
Differences of branch points are killed by $31$, and $c$
permutes those branch points; hence $cD_0-D_0$ is torsion.
For the quotient map $\pi:C\to B$,
$\pi^*\pi_*D_0=(1+c+c^2)D_0$ is $3D_0$ plus torsion.
Thus $\pi_*D_0\in A(\Q)$ has infinite order.

For the upper bound use the equivalent triple $(5,1,-6)$.
The numerators of $B_2,B_4,\ldots,B_{28}$ have respective
nonzero residues modulo $31$
\[
 1,30,1,30,5,22,7,10,2,12,29,28,17,28.
\]
Kummer's criterion makes $31$ regular. If
$q(u)=(u^{30}-1)/31\pmod{31}$, direct integer arithmetic gives
\[
 q(5^5(-6)^{-6})=0,\qquad B_{28}=14\pmod{31},\qquad
 q(5^5(-6)^{-6})^3-30B_{28}=14\pmod{31}.
\]
The quotient is therefore tame. The required odd unit eigenspace
in \citet[Theorem~1.2 and the following criterion]{McCallumTzermias2004}
vanishes because $31\nmid B_{14}$.
\citet[Theorems~4.16--4.17 and Corollary~4.18, author
pp.~35--36]{ShuExpanded2021} now gives
$\dim_{\mathbb F_{31}}\Sel_\lambda(J/K)=7$.
Here $\lambda=\zeta_{31}-1$; replacing it by
$\zeta_{31}-\zeta_{31}^{-1}$ changes it by a unit.
The tame parameter is zero, for which that source explicitly uses
$(0/31)=-1$ in its dimension formula.

Write $R=\operatorname{rank}J(\Q)$. Semilinear Galois descent
on the $K$-vector space $J(K)\otimes\Q$ gives dimension $R$
over $K$. Its finite torsion contributes one dimension modulo
$\lambda$, since $J[\lambda]$ has order $31$ and is rational
over $K$. Thus the Selmer exact sequence gives
\[
 7=R+1+\dim_{\mathbb F_{31}}\Sha(J/K)[\lambda].
\]
\citet[Theorems~1.2 and 1.4; Proposition~5.1]{McCallumTzermias2004}
proves that $\Sha[\lambda^3]$ is free over
$\mathcal O_K/(\lambda^3)$ and that its pairing with
$\Sha[\lambda]$ is nonzero. The latter group cannot have
dimension one: if $a$ generated the former, $\lambda^2a$
would generate the latter, whereas the Rosati adjoint satisfies
$\widehat\lambda^2\equiv\lambda^2\pmod{\lambda^3}$.
Functoriality and alternation would then give
\[
 [a,\lambda^2a]=[\widehat\lambda^2a,a]
              =[\lambda^2a,a]=-[a,\lambda^2a],
\]
forcing the pairing to vanish at the odd prime $31$.
Consequently $R\le4$. Since $A(\Q)$ has positive rank,
$\operatorname{rank}P(\Q)\le3$; its $\Q(\zeta_3)$-action
makes that rank even, proving the assertion.
\end{proof}

\begin{proposition}[The exact CM class numbers]
\label{prop:thirtyone-cm-class-numbers}
The fields $F_0$ and $M$ have discriminants $31^4$ and $-31^9$,
and class numbers
\[
 h(F_0)=1,\qquad h(M)=3.
\]
The type norm for $T=\{0,1,2,4,8\}$ acts trivially on
$\operatorname{Cl}(M)$.
\end{proposition}
\begin{proof}
The conductor--discriminant formula gives the displayed
discriminants, since every nontrivial character has conductor $31$.
For $R_0=\{\pm1,\pm5,\pm25\}$, put
$\eta=\sum_{a\in R_0}\zeta_{31}^a$.
In the basis of the five Gaussian periods indexed by
$3^jR_0$, cyclic convolution gives the multiplication matrix
\[
 \begin{pmatrix}
 -4&0&0&2&1\\ -6&1&2&1&2\\ -6&2&2&1&1\\
 -4&1&1&0&2\\ -5&2&1&2&0
 \end{pmatrix}.
\]
Its characteristic polynomial is
$t^5+t^4-12t^3-21t^2+t+5$, so $N(\eta)=-5$.
The Minkowski bound is
$5!\,31^2/5^5=115320/3125<37$.
The only prime ideals below that bound lie over $2,5,31$:
the first has norm $2^5=32$ and is the principal ideal $(2)$;
the five primes over $5$ are generated by the conjugates of
$\eta$; the prime over $31$ is generated by
$N_{K/F_0}(1-\zeta_{31})$. This exhausts the bounded ideals,
and proves $h(F_0)=1$.

The roots of unity in $M$ are $\{\pm1\}$. For a unit $u$,
Kronecker's theorem makes $u/\bar u$ a root of unity.
If it were $-1$, then $M=F_0(u)$ with $u^2$ a negative unit.
That quadratic extension would have discriminant dividing $4$,
contrary to ramification at $31$. Thus every unit is real and
the Hasse unit index is one. Let $z$ be a primitive tenth root
of unity and use the character taking $3\bmod31$ to $z$.
Reduction modulo $\Phi_{10}(z)$ gives the exact identity
\[
 \sum_{j=0}^{29}(3^j\bmod31)z^j=-62z.
\]
Hence the four primitive odd order-ten characters have
Bernoulli product $16$. The remaining odd quadratic character has
\[
 B_{1,\chi}=\frac1{31}\sum_{j=0}^{29}(3^j\bmod31)(-1)^j=-3.
\]
The analytic class-number formula, applied to the Dirichlet
factorization of $\zeta_M/\zeta_{F_0}$, therefore gives
\[
 \frac{h(M)}{h(F_0)}=\frac{2}{2^5}
       \prod_{\chi\ \mathrm{odd}}(-B_{1,\chi})
       =\frac{3\cdot16}{16}=3.
\]
Finally complex conjugation acts by inversion on $\operatorname{Cl}(M)$,
since $h(F_0)=1$. A generator $\sigma$ of its cyclic Galois
group of order ten consequently acts as $-1$ on this cyclic
class group of order three, because $\sigma^5$ is conjugation.
The type has four even and one odd exponent, so its norm acts
by $4-1=3$, and is zero on the class group.
\end{proof}

\begin{proposition}[The central vanishing and the cubic-character restriction]
\label{prop:thirtyone-central-obstruction}
The whole-Jacobian function has root number $+1$ and
\[
 L(J/\Q,1)=L'(J/\Q,1)=L(A/\Q,1)=0.
\]
If $\chi$ is a cubic Dirichlet character modulo $31$, then,
over a coefficient field containing its values,
\[
 V(J)=V(A)\oplus V(A)\otimes\chi\oplus V(A)\otimes\chi^2.
\]
The two complementary weight-two Hilbert constituents over $F_0$
have nontrivial cubic central characters. Neither the rank-zero
assertion for $P$ nor finiteness of its Tate--Shafarevich group
is established by this calculation.
\end{proposition}
\begin{proof}
The restriction of $\chi$ to $H$ is nontrivial, since
$5=3^{20}\pmod{31}$. If $\varphi$ is the algebraic CM Hecke
character of $M$ associated with $A$, and
$\nu=\chi\circ N_{M/\Q}$, then $\nu$ cuts out $K/M$.
The three characters of this cyclic cubic extension give
\[
 L(J/\Q,s)=L_M(\varphi,s)L_M(\varphi\nu,s)
                            L_M(\varphi\nu^2,s),\qquad
 L(A/\Q,s)=L_M(\varphi,s).
\]
Indeed the three-dimensional multiplicity space over $G_M$
has these three characters, and $c$-invariants select the
untwisted constituent. Inducing to $G_\Q$ proves the asserted
representation identity.

Quadratic induction from $M/F_0$ gives a weight-two Hilbert
form $\pi$. Its normalized central character is trivial:
real multiplication on $A_{F_0}$ and its polarization give
cyclotomic determinant on every two-dimensional constituent.
As in Proposition~\ref{prop:lm-nineteen-rankone},
the CM construction, the attached-variety statement and Faltings'
isogeny theorem identify the abelian variety to which
\citet[Theorem~A, p.~28]{Zhang2001} applies.
The real degrees are $[F_0:\Q]=5$ and $[K^+:\Q]=15$, both
odd. A nonzero central value for $A/\Q$, respectively $J/\Q$,
would therefore force rank zero over these real fields.
The rational infinite-order points already proved rule this
out. The tame root-number formula
of \citet[Theorem~3.13 and the tame case in \S5]{ShuExpanded2021}
gives $W(J/\Q)=(2/31)=+1$.
Its functional equation then makes the central order even,
so $L'(J/\Q,1)=0$ as well.

Put $\chi_{F_0}=\chi\circ N_{F_0/\Q}$.
The two complementary forms $\pi\otimes\chi_{F_0}$ and
$\pi\otimes\chi_{F_0}^2$ have central characters
$\chi_{F_0}^2$ and $\chi_{F_0}$, respectively.
They are nontrivial because $[F_0:\Q]=5$ and the character
cuts out a cubic extension. The trivial-central-character
hypothesis in Zhang's theorem therefore excludes these direct
applications. Adjoining the real cubic field kills the twists,
but makes $P$ isogenous to two copies of the positive-rank $A$
after base change; it does not produce a rank-zero application.
\end{proof}

For comparison with another possible arithmetic input, use the
September 2025 author version of \citet{Hsieh2014}, whose
Theorem~2 and Corollaries~2--3 are on pp.~3--4.
Its inclusion $(F_{\psi,\Sigma})\subset(L_{\psi,\Sigma})$
is a lower-bound divisibility for the Selmer module, so it does
not bound that module above from a nonzero central value.
Corollary~2 supplies equality for an anticyclotomic finite branch,
but that condition fails here already at $31$.
Indeed $\varphi^c\varphi$ is the norm character and
$\nu^c=\nu$, so
$(\varphi\nu)^c(\varphi\nu)=\operatorname{Norm}\,\nu^2$.
For any permitted coefficient prime different from $31$ and $3$,
the cyclotomic character is unramified at $31$, whereas $\nu^2$
has nontrivial tame cubic inertia. This obstruction survives
passage to the finite branch and dual or Tate twists.
Corollary~3 instead requires a CM type induced from an imaginary
quadratic field. Although $M$ contains $\Q(\sqrt{-31})$, the
primitive type $T$ contains both parities and is not so induced.
These are failures of the specified hypotheses, not a claim that
no other arithmetic theorem could apply.

The nonvanishing of $L(P/\Q,1)$ remains unproved here.
Even such a certificate would still require an applicable Selmer
upper-bound or descent theorem for the cubic complementary
constituents. The present rank bound leaves
$\operatorname{rank}P(\Q)=0$ or $2$; no exponent-thirty-one
finite-descent or section conclusion is asserted.
